%% LyX 2.4.4 created this file.  For more info, see https://www.lyx.org/.
%% Do not edit unless you really know what you are doing.
\documentclass[english]{article}
\usepackage[T1]{fontenc}
\usepackage[latin9]{inputenc}
\usepackage{units}
\usepackage{enumitem}
\usepackage{amsmath}
\usepackage{amsthm}
\usepackage{amssymb}
\usepackage{geometry}
\geometry{verbose,tmargin=1in,bmargin=1in,lmargin=1in,rmargin=1in}

\makeatletter
%%%%%%%%%%%%%%%%%%%%%%%%%%%%%% Textclass specific LaTeX commands.
\numberwithin{equation}{section}
\numberwithin{figure}{section}
\newlength{\lyxlabelwidth}      % auxiliary length 

%%%%%%%%%%%%%%%%%%%%%%%%%%%%%% User specified LaTeX commands.
\usepackage{fancyhdr}
\usepackage{lastpage}
\pagestyle{fancy}
\usepackage{enumitem}
\usepackage{ifthen}
\everymath=\expandafter{\the\everymath\displaystyle}
%\DeclareMathSizes{10}{10}{10}{10}
\usepackage{multicol}

\usepackage{tikz}
\usepackage{tikz-3dplot}
\usetikzlibrary{calc,arrows}

\usetikzlibrary{patterns}
\usetikzlibrary{plotmarks}
\usepackage{pgfplots}
\pgfplotsset{compat=newest}

\usepackage{calc}
\usepackage[nomessages]{fp}

\usepackage{datetime}
% Added by lyx2lyx
\setlength{\parskip}{\smallskipamount}
\setlength{\parindent}{0pt}

\makeatother

\usepackage{babel}
\begin{document}
\renewcommand{\author}{Dr.\,James Ashe}

\newcommand{\classNum}{336}

\newcommand{\classSec}{A}

\newcommand{\examType}{Exam 2}

\newcommand{\nameType}{Solutions}

\renewcommand{\date}{\formatdate{25}{10}{2013}}

%%First page's header%%%%%%%%%%%%%%%%%%%%%%
\fancypagestyle{firststyle} %%header settings for first pagr
{
	\fancyhead[L]{MTH \classNum{}(\classSec{}) \author{}\\
	\textbf{\examType{}} ~\date{}}
	\fancyhead[C]{\textbf{\nameType}}
	\fancyhead[R]{}
	\setlength{\headsep}{14pt}
}
%%%%%%%%%%%%%%%%%%%%%%%%%%%%%%%%%%%%%%%%%%

%%Next page's header%%%%%%%%%%%%%%%%%%%%%%%
\setlist{nolistsep}
\lhead{MTH \classNum{}(\classSec{})} 
\chead{\textbf{\nameType}} 
\cfoot{\thepage{}~of~\pageref{LastPage}} 
\setlength{\headsep}{5pt} 
%%%%%%%%%%%%%%%%%%%%%%%%%%%%%%%%%%%%%%%%%%%

%%Various settings; see the preamble for other settings%%%%%
%\setenumerate[0]{leftmargin=12pt,itemindent=0pt,labelwidth=0pt}
\setlist{topsep=.5pt}
\renewcommand{\labelenumi}{\textbf{\arabic{enumi}.}}
\renewcommand{\labelenumii}{\textbf{\alph{enumii}.}}
\thispagestyle{firststyle}
%%%%%%%%%%%%%%%%%%%%%%%%%%%%%%%%%%%%%%%%%%

\newcommand{\xmax}{0}
\newcommand{\xmin}{-\xmax}
\newcommand{\xstep}{0}
\newcommand{\ymax}{0}
\newcommand{\ymin}{-\ymax}
\newcommand{\ystep}{0} 
\newcommand{\dspace}{\vspace{1cm}}
\newcommand{\xa}{0}
\newcommand{\xb}{0}
\newcommand{\ya}{1}
\newcommand{\yb}{1}

\newcommand{\xspace}{\vspace{0cm}}
\newcommand{\Xspace}{\vspace{0cm}}

\newcommand{\vectors}[2]{\textbf{#1}_{1},\textbf{#1}_{2},\dots,\textbf{#1}_{#2}}
\newcommand{\vectorsN}[1]{\textbf{#1}_{1},\textbf{#1}_{2},\dots,\textbf{#1}_{n}}
\newcommand{\R}[1]{\mathbb{R}^{#1}}
\begin{enumerate}[resume]
\item Let $A=\begin{bmatrix}\begin{array}{rr}
1 & 2\\
3 & 4
\end{array}\end{bmatrix}$, $B=\begin{bmatrix}\begin{array}{rr}
-1 & 1\\
0 & 1
\end{array}\end{bmatrix}$, and $C=\begin{bmatrix}\begin{array}{rrr}
1 & 2 & 2\\
0 & 2 & 4
\end{array}\end{bmatrix}$. Calculate each of the following expressions or explain why it is
not defined.

\begin{enumerate}[resume]
\item $AB$

\begin{enumerate}[resume]
\item []$AB=\begin{bmatrix}\begin{array}{rr}
1(-1)+2(0) & 1(1)+2(1)\\
3(-1)+4(0) & 3(1)+4(1)
\end{array}\end{bmatrix}=\begin{bmatrix}\begin{array}{rr}
-1 & 3\\
-3 & 7
\end{array}\end{bmatrix}$\Xspace \vfill
\end{enumerate}
\item $CB$

\begin{enumerate}[resume]
\item []$CB$ is undefined because the $C$ has 3 columns but $B$ has
two rows.\xspace \vfill
\end{enumerate}
\end{enumerate}
\item Suppose that $T\colon\R{3}\rightarrow\R{2}$ is defined by $T\left(\begin{bmatrix}\begin{array}{r}
x_{1}\\
x_{2}\\
x_{3}
\end{array}\end{bmatrix}\right)=\begin{bmatrix}\begin{array}{r}
2x_{1}\\
x_{2}+x_{3}
\end{array}\end{bmatrix}$.

\begin{enumerate}
\item Find the standard matrix $A$ for $T$.

\begin{enumerate}
\item []$T\left(\begin{bmatrix}\begin{array}{r}
1\\
0\\
0
\end{array}\end{bmatrix}\right)=\begin{bmatrix}\begin{array}{r}
2\\
0
\end{array}\end{bmatrix}$, $T\left(\begin{bmatrix}\begin{array}{r}
0\\
1\\
0
\end{array}\end{bmatrix}\right)=\begin{bmatrix}\begin{array}{r}
0\\
1
\end{array}\end{bmatrix}$, and $T\left(\begin{bmatrix}\begin{array}{r}
0\\
0\\
1
\end{array}\end{bmatrix}\right)=\begin{bmatrix}\begin{array}{r}
0\\
1
\end{array}\end{bmatrix}$. So $A=\begin{bmatrix}\begin{array}{rrr}
2 & 0 & 0\\
0 & 1 & 1
\end{array}\end{bmatrix}$.\Xspace \vfill
\end{enumerate}
\end{enumerate}
\begin{enumerate}[resume]
\item Is $T$ is a linear transformation? 

\begin{enumerate}[resume]
\item []Yes. You can check the properties or see that $\begin{bmatrix}\begin{array}{rrr}
2 & 0 & 0\\
0 & 1 & 1
\end{array}\end{bmatrix}\begin{bmatrix}\begin{array}{r}
x_{1}\\
x_{2}\\
x_{3}
\end{array}\end{bmatrix}=\begin{bmatrix}\begin{array}{r}
2x_{1}\\
x_{2}+x_{3}
\end{array}\end{bmatrix}$.\xspace vfill
\end{enumerate}
\item Let $\mathbf{x}\in\R{n}$ and $A$ be an $m\times n$ matrix. The
function $f(\mathbf{x})=A\mathbf{x}$ is a linear map from \underline{\hspace{.15cm}$\R{m}$\hspace{.15cm}}
to \underline{\hspace{.15cm}$\R{n}$\hspace{.15cm}}.\xspace \vfill
\end{enumerate}
\item Find $A^{-1}$ or show that it does not exist.

\begin{enumerate}
\item $A=\begin{bmatrix}1 & 1 & 0\\
0 & 1 & 1\\
0 & 0 & 1
\end{bmatrix}$

\begin{enumerate}
\item []$\begin{bmatrix}\begin{array}{rrr|rrr}
1 & 1 & 0 & 1 & 0 & 0\\
0 & 1 & 1 & 0 & 1 & 0\\
0 & 0 & 1 & 0 & 0 & 1
\end{array}\end{bmatrix}\rightsquigarrow\begin{array}{r}
\\\mbox{R2-R3}\\
\\\end{array}\begin{bmatrix}\begin{array}{rrr|rrr}
1 & 1 & 0 & 1 & 0 & 0\\
0 & 1 & 0 & 0 & 1 & -1\\
0 & 0 & 1 & 0 & 0 & 1
\end{array}\end{bmatrix}\rightsquigarrow\begin{array}{r}
\mbox{R1-R2}\\
\\\\\end{array}\begin{bmatrix}\begin{array}{rrr|rrr}
1 & 0 & 0 & 1 & -1 & 0\\
0 & 1 & 0 & 0 & 1 & -1\\
0 & 0 & 1 & 0 & 0 & 1
\end{array}\end{bmatrix}$ so $A^{-1}=\begin{bmatrix}\begin{array}{rrr}
1 & -1 & 0\\
0 & 1 & -1\\
0 & 0 & 1
\end{array}\end{bmatrix}$.\Xspace \vfill
\end{enumerate}
\end{enumerate}
\begin{enumerate}[resume]
\item $A=\begin{bmatrix}2 & 2 & 0\\
0 & 2 & 2\\
0 & 0 & 0
\end{bmatrix}$

\begin{enumerate}[resume]
\item []$A^{-1}$ does not exist since $A$ is singular.\xspace \vfill
\end{enumerate}
\end{enumerate}
\end{enumerate}
\newpage
\begin{enumerate}[resume]
\item Let $A=\begin{bmatrix}\begin{array}{rr}
\nicefrac{1}{2} & \nicefrac{1}{4}\\
0 & \nicefrac{1}{2}
\end{array}\end{bmatrix}$, $A^{-1}=\begin{bmatrix}\begin{array}{rr}
2 & -1\\
0 & 2
\end{array}\end{bmatrix}$ and $\mathbf{b}=\begin{bmatrix}\begin{array}{r}
1\\
1
\end{array}\end{bmatrix}$

\begin{enumerate}
\item Use $A^{-1}$ to solve $A\mathbf{x}=\mathbf{b}$ for \textbf{$\mathbf{x}$}.

\begin{enumerate}
\item []$A\mathbf{x}=\mathbf{b}\Rightarrow\mathbf{x}=A^{-1}\mathbf{b}$
so $\mathbf{x}=\begin{bmatrix}\begin{array}{rr}
2 & -1\\
0 & 2
\end{array}\end{bmatrix}\begin{bmatrix}\begin{array}{r}
1\\
1
\end{array}\end{bmatrix}=\begin{bmatrix}\begin{array}{r}
1\\
2
\end{array}\end{bmatrix}$\Xspace \vfill
\end{enumerate}
\end{enumerate}
\begin{enumerate}[resume]
\item Use your answer from part \textbf{(a)} to express $\mathbf{b}$ as
a linear combination of $A$.

\begin{enumerate}[resume]
\item []$1(\nicefrac{1}{2},0)+2(\nicefrac{1}{4},\nicefrac{1}{2})=(1,1)$\Xspace \vfill
\end{enumerate}
\end{enumerate}
\item Let $\mathbf{v}_{1}=(1,2)$, $\mathbf{v}_{2}=(0,1)$, and $\mathbf{v}_{3}=(1,0)$
be three vectors in $\R{2}$.

\begin{enumerate}
\item Show that $\mathbf{v}_{1}=(1,2)$ and $\mathbf{v}_{2}=(0,1)$ gives
a basis for $\R{2}$.

\begin{enumerate}
\item []$\begin{bmatrix}\begin{array}{rr}
1 & 0\\
2 & 1
\end{array}\end{bmatrix}\rightsquigarrow\begin{array}{r}
\\\mbox{R2-2R1}
\end{array}\begin{bmatrix}\begin{array}{rr}
1 & 0\\
0 & 1
\end{array}\end{bmatrix}$ which is nonsingular so it is a basis.\Xspace \vfill
\end{enumerate}
\end{enumerate}
\begin{enumerate}[resume]
\item Find another basis for $\R{2}$.

\begin{enumerate}[resume]
\item []There are many answers; the standard basis is one of them, $\{(1,0),(0,1)\}$.\vfill
\end{enumerate}
\item Is $\mathbf{v}_{1}=(1,2)$, $\mathbf{v}_{2}=(0,1)$, and $\mathbf{v}_{3}=(1,0)$
linearly independent? Explain.

\begin{enumerate}[resume]
\item []No. The $\dim\R{2}=2$ and this set has three vectors.\vfill
\end{enumerate}
\item Does $\mathbf{v}_{1}=(1,2)$, $\mathbf{v}_{2}=(0,1)$, and $\mathbf{v}_{3}=(1,0)$
span $\R{2}$?

\begin{enumerate}[resume]
\item []Yes. Since $\mbox{Span}\{(1,2),(0,1)\}=\R{2}$, $(1,0)\in\mbox{Span}\{(1,2),(0,1)\}\Rightarrow\mbox{Span}\{(1,2),(0,1),(1,0\}=\mbox{Span}\{(1,2),(0,1)\}=\R{2}$.
\vfill\newpage
\end{enumerate}
\end{enumerate}
\item Let $\mathbf{v}_{1}=(2,4,8)$, $\mathbf{v}_{2}=(-2,-4-8)$, $\mathbf{v}_{3}=(1,2,4)$
be vectors in $\R{3}$. 

\begin{enumerate}
\item Is $\{\mathbf{v}_{1},\mathbf{v}_{2},\mathbf{v}_{3}\}$ linearly independent?

\begin{enumerate}
\item []No. $(2,4,8)=-1(-2,-4,-8)\Rightarrow(2,4,8)-(-2,-4,-8)+0(1,2,4)=\mathbf{0}$.
You can also find the reduced echelon form the asscoiated column matrix
and see that it is singular (rank of 1).\\
$A=\begin{bmatrix}\begin{array}{rrr}
2 & -2 & 1\\
4 & -4 & 2\\
8 & -8 & 4
\end{array}\end{bmatrix}\rightsquigarrow\begin{bmatrix}\begin{array}{rrr}
2 & -2 & 1\\
0 & 0 & 0\\
0 & 0 & 0
\end{array}\end{bmatrix}$\xspace\Xspace \vfill
\end{enumerate}
\end{enumerate}
\begin{enumerate}[resume]
\item Does $\{\mathbf{v}_{1},\mathbf{v}_{2},\mathbf{v}_{3}\}$ span $\R{3}$?,
i.e., does $\mbox{Span}(\mathbf{v}_{1},\mathbf{v}_{2},\mathbf{v}_{3})=\R{3}$?
Explain.

\begin{enumerate}[resume]
\item []No. The $\mbox{dim}\R{3}=3$ but $\{\mathbf{v}_{1},\mathbf{v}_{2},\mathbf{v}_{3}\}$
is linearly dependent (meaning that a basis will have less than 3
vectors).\vfill
\end{enumerate}
\item Find the dimension of $\mbox{Span}(\mathbf{v}_{1},\mathbf{v}_{2},\mathbf{v}_{3})$.

\begin{enumerate}[resume]
\item []From part \textbf{(a) }we can see that $\mbox{rank}A=1$, so a
basis will have 1 vector and the dimension must be 1.\vfill
\end{enumerate}
\item Find a basis for $\mbox{Span}(\mathbf{v}_{1},\mathbf{v}_{2},\mathbf{v}_{3})$.

\begin{enumerate}[resume]
\item []From part \textbf{(a) }we can see that $(2,4,8)$ is a basis.\vfill
\end{enumerate}
\end{enumerate}
\item Let $V=\{\vectors{v}{n+1}\}$ be a set of $n+1$ vectors in $\mathbb{R}^{n}$.

\begin{enumerate}
\item Is $V$ linearly independent or linearly dependent? Explain.

\begin{enumerate}
\item []It is linearly dependent. A basis for $\R{n}$ will have $n$ vectors
($\dim\R{n}=n$). So then a set of $n+1$ vectors must be linearly
dependent.\vfill
\end{enumerate}
\end{enumerate}
\begin{enumerate}[resume]
\item Is $V$ a basis for $\mathbb{R}^{n}$? Explain. 

\begin{enumerate}[resume]
\item []No. It is linearly dependent.\vfill
\end{enumerate}
\item Is $\mbox{Span}V$ a subspace of $\mathbb{R}^{n}$?

\begin{enumerate}[resume]
\item []Yes a span is always a subspace.\vfill
\end{enumerate}
\end{enumerate}
\item [\textbf{Bonus.}]Let $\mathbf{v}_{1},\mathbf{v}_{2},\dots,\mathbf{v}_{k}\in\mathbb{R}^{n}$
and let $\mathbf{v}\in\mathbb{R}^{n}$. Prove that $\mathbf{v}\in\mbox{Span}\left(\mathbf{v}_{1},\mathbf{v}_{2},\dots,\mathbf{v}_{k}\right)$
$\iff$ $\mbox{Span}\left(\mathbf{v}_{1},\mathbf{v}_{2},\dots,\mathbf{v}_{k}\right)=\mbox{Span}\left(\mathbf{v}_{1},\mathbf{v}_{2},\dots,\mathbf{v}_{k},\mathbf{v}\right)$.

\begin{proof}
$(\Rightarrow)$ Let $\mathbf{v}\in\mbox{Span}\left(\mathbf{v}_{1},\mathbf{v}_{2},\dots,\mathbf{v}_{k}\right)$.
Then by definition there are real numbers $c_{1},c_{2},\dots,c_{k}$
such that $\mathbf{v}=c_{1}\mathbf{v}_{1}+c_{2}\mathbf{v}_{2}+\dots+c_{k}\mathbf{v}_{k}$
so then $\mbox{Span}\left(\mathbf{v}_{1},\mathbf{v}_{2},\dots,\mathbf{v}_{k},\mathbf{v}\right)\subset\mbox{Span}\left(\mathbf{v}_{1},\mathbf{v}_{2},\dots,\mathbf{v}_{k}\right)$.
Also $\mbox{Span}\left(\mathbf{v}_{1},\mathbf{v}_{2},\dots,\mathbf{v}_{k}\right)\subset\mbox{Span}\left(\mathbf{v}_{1},\mathbf{v}_{2},\dots,\mathbf{v}_{k},\mathbf{v}\right)$
so then $\mbox{Span}\left(\mathbf{v}_{1},\mathbf{v}_{2},\dots,\mathbf{v}_{k}\right)=\mbox{Span}\left(\mathbf{v}_{1},\mathbf{v}_{2},\dots,\mathbf{v}_{k},\mathbf{v}\right)$.

$(\Leftarrow)$ Let $\mbox{Span}\left(\mathbf{v}_{1},\mathbf{v}_{2},\dots,\mathbf{v}_{k}\right)=\mbox{Span}\left(\mathbf{v}_{1},\mathbf{v}_{2},\dots,\mathbf{v}_{k},\mathbf{v}\right)$.
$\mathbf{v}\in\mbox{Span}\left(\mathbf{v}_{1},\mathbf{v}_{2},\dots,\mathbf{v}_{k},\mathbf{v}\right)$
so then by the hypothesis $\mathbf{v}\in\mbox{Span}\left(\mathbf{v}_{1},\mathbf{v}_{2},\dots,\mathbf{v}_{k}\right)$.
\end{proof}
\end{enumerate}

\end{document}
